\documentclass[11pt,letterpaper]{article}
\usepackage[margin=1in]{geometry}
\usepackage[T1]{fontenc}
\usepackage{lmodern}
\usepackage{amsmath,amssymb,amsthm}
\usepackage{microtype}
\usepackage{booktabs,longtable,array}
\usepackage{enumitem}
\usepackage{xcolor}
\usepackage{xurl}
\usepackage{hyperref}
\usepackage{fancyhdr}
\definecolor{inkblue}{RGB}{27,55,78}
\hypersetup{colorlinks=true,linkcolor=inkblue,citecolor=inkblue,urlcolor=inkblue,
 pdftitle={Smaller explicit Diophantine forms of ZFC consistency},
 pdfauthor={AI-assisted research for the ProveIt repository}}
\newcommand{\Z}{\mathbb Z}
\newcommand{\N}{\mathbb N}
\DeclareMathOperator{\Pos}{Pos}
\newtheorem{theorem}{Theorem}[section]
\newtheorem{lemma}[theorem]{Lemma}
\newtheorem{proposition}[theorem]{Proposition}
\theoremstyle{remark}
\newtheorem{remark}[theorem]{Remark}
\setcounter{tocdepth}{2}
\setlength{\headheight}{14pt}
\setlength{\emergencystretch}{2em}
\pagestyle{fancy}
\fancyhf{}
\fancyhead[L]{\small\textsc{Diophantine forms of ZFC consistency}}
\fancyhead[R]{\small ProveIt synthesis}
\fancyfoot[C]{\thepage}
\renewcommand{\headrulewidth}{0.3pt}
\setlist{itemsep=3pt,topsep=5pt}
\allowdisplaybreaks[1]

\begin{document}
\hypersetup{pageanchor=false}
\begin{titlepage}
\vspace*{0.55in}
{\color{inkblue}\Large\textsc{ProveIt research report}\par}
\vspace{0.35in}
{\huge\bfseries Smaller explicit Diophantine\par
 forms of ZFC consistency\par}
\vspace{0.3in}
{\Large Twenty-four integers, tighter twelve-variable norms,\par
 and smaller materialized coefficients\par}
\vspace{0.35in}
{\large AI-assisted mathematical synthesis\par}
{\large Prepared for Vladimir Reshetnikov\par}
\vspace{0.15in}
{22 October 2026\par}
\vspace{0.4in}
\noindent\rule{\textwidth}{0.6pt}
\begin{abstract}
\noindent
Combining rational Pell saturation, separating-root norms, signed-domain
arguments and a concrete padding controller gives an explicit polynomial
in 24 unrestricted integers of degree 1936 with a 167-operation circuit.
ZFC proves that its lack of integer solutions is equivalent to
$\operatorname{Con}(\mathrm{ZFC})$.  Its three-equation version has degree
1920 and uses 162 operations.  Three twelve-positive-variable alternatives
have degree/operation pairs $(1069,123)$, $(1073,121)$ and $(1081,120)$.
An exact refactoring reduces the existing degree-1353 twelve-variable
equation to 113 operations.  Retaining the 943-tile language while
replacing five-power rounding reduces the fixed base to
$B=2^{6364399}$; all four resulting integers are materialized.
The report supplies explicit polynomial definitions, ordinary
mathematical proofs, arithmetic circuits and reproducible finite audits.
It does not claim a global minimum or a new kernel-checked end-to-end
formalization.
\end{abstract}
\vfill
\noindent\small
Source snapshot: \texttt{ae3007ca4467011e0e051c84acb4a1dd13971b0e}.\\
The archive preserves source provenance and separates new proofs from
inherited compiler and proof-calibration results.
\end{titlepage}

\hypersetup{pageanchor=true}
\pagenumbering{roman}
\tableofcontents
\clearpage
\pagenumbering{arabic}
\input{sections/front.tex}
\input{sections/positive.tex}
\input{sections/integer.tex}
\input{sections/circuits.tex}
\input{sections/coefficients.tex}
\input{sections/reproduction.tex}
\input{sections/references.tex}
\end{document}
